\begin{abstractpage}

\begin{abstract}{romanian}
Această lucrare investighează în ce măsură modul în care un model de detecție a anomaliilor reprezintă structura latentă a datelor \enquote{normale} determină tipurile de anomalii pe care le poate detecta. Șapte modele, acoperind patru paradigme structural distincte (graniță, partiționare, bottleneck de compresie, obiectiv predictiv), sunt evaluate printr-un pipeline uniform pe trei seturi de date cu grade diferite de structură temporală -- tranzacții bancare tabelare, serii de timp Yahoo Webscope S5 și fluxuri de rețea CIC-IDS2017.
Rezultatele arată că performanța fiecărei paradigme depinde sistematic de tipul de anomalie și de structura datelor, nu de puterea intrinsecă a modelului, iar o serie de studii de caz de explicabilitate confirmă direct mecanismele propuse pentru a explica aceste diferențe.
Un ansamblu construit pe baza a două paradigme cu evidențe demonstrabil disjuncte extinde acoperirea pe clase de atac altfel nedetectabile individual, confirmând valoarea combinării unor asumpții complementare despre normalitate.
\end{abstract}

\begin{abstract}{english}
This thesis investigates the extent to which the way an anomaly detection model represents the latent structure of \enquote{normal} data determines the types of anomalies it can detect. Seven models, spanning four structurally distinct paradigms (boundary, partitioning, compression bottleneck, predictive objective), are evaluated under a uniform pipeline across three datasets with varying degrees of temporal structure -- tabular credit card transactions, Yahoo Webscope S5 time series, and CIC-IDS2017 network flows.
Results show that each paradigm's performance depends systematically on the type of anomaly and the structure of the data, rather than on the model's intrinsic strength, and a series of explainability case studies directly confirm the mechanisms proposed to explain these differences.
An ensemble built from two paradigms with demonstrably disjoint evidence extends coverage to attack classes otherwise undetectable individually, confirming the value of combining complementary assumptions about normality.
\end{abstract}

\end{abstractpage}